% Part: second-order-logic
% Chapter: syntax-and-semantics
% Section: introduction

\documentclass[../../../include/open-logic-section]{subfiles}

\begin{document}

\olfileid{sol}{syn}{int}

\olsection{Pambuka}

Ing logika orde kapisan, kita nggabungake simbol nonlogis saka
sawijining basa, yaiku !!{constant}s, !!{function}s, lan
!!{predicate}s, karo simbol logis kanggo nyatakake perkara ngenani
!!{structure}s orde kapisan. Iki ditindakake nganggo gagasan
relasi nyukupi, kang ngubungake !!a{structure}~$\Struct{M}$,
sesarengan karo pangaji variabel~$s$, lan !!a{formula}~$!A$:
$\Sat{M}{!A}[s]$ bener yen lan mung yen apa kang dinyatakake
dening $!A$ iku bener nalika !!{constant}s, !!{function}s, lan
!!{predicate}s diinterpretasikake miturut $\Struct{M}$, sarta
variabel bebase diinterpretasikake miturut~$s$. Interpretasi
!!{identity}~$\eq$ wis kalebu ing definisi
$\Sat{M}{!A}[s]$, semono uga interpretasi~$\lforall$ lan
$\lexists$. Kang kapisan tansah diinterpretasikake minangka relasi
identitas ing !!{domain}~$\Domain{M}$ saka struktur kasebut,
dene kuantor tansah diinterpretasikake minangka nyakup sakabehe
!!{domain}. Nanging kang wigati, kuantifikasi mung kena tumrap
unsur-unsur saka !!{domain}, mula mung !!{variable}s objek kang
kena mapan sawisé kuantor.

Ing logika orde kapindho, basa lan definisi relasi nyukupi
dijembarake supaya ngemot variabel fungsi lan predikat kang bebas
utawa kaiket, uga kuantifikasi tumrap variabel kasebut.
Sesambungane variabel-variabel iki karo !!{function}s lan
!!{predicate}s padha kaya sesambungane variabel objek karo
!!{constant}s. Kalungguhane ing pambentukan term lan !!{formula}s
logika orde kapindho uga padha, lan kuantifikasi tumrape ditangani
kanthi cara kang saemper. Ing semantik \emph{standar}, kuantor
orde kapindho nyakup sakabehe objek kang jinise cocog:
fungsi $n$-papan saka $\Domain{M}$ menyang~$\Domain{M}$ kanggo
variabel fungsi, lan relasi $n$-papan kanggo variabel predikat.
Tuladhane, sanadyan
$\lforall[\Obj{v_0}][(\Obj{P^1_0}(\Obj{v_0}) \lor \lnot
  \Obj{P^1_0}(\Obj{v_0}))]$ iku formula ing logika orde kapisan
lan orde kapindho, ing kang kapindho kita uga bisa nimbang
$\lforall[\Obj{V^1_0}][\lforall[\Obj{v_0}][(\Obj{V^1_0}(\Obj{v_0}) \lor
    \lnot \Obj{V^1_0}(\Obj{v_0}))]]$ lan
$\lexists[\Obj{V^1_0}][\lforall[\Obj{v_0}][(\Obj{V^1_0}(\Obj{v_0}) \lor
    \lnot \Obj{V^1_0}(\Obj{v_0}))]]$. Amarga loro-lorone ora ngemot
variabel bebas, loro-lorone iku !!{sentence}s logika orde kapindho.
Ing kene, $\Obj{V^1_0}$ iku variabel predikat $1$-papan orde
kapindho. Interpretasi kang diidinake kanggo $\Obj{V^1_0}$
padha karo kang bisa kita wenehake marang !!{predicate} $1$-papan
kaya $\Obj{P^1_0}$, yaiku subhimpunan
saka~$\Domain{M}$. Dadi, kuantifikasi tumrape ateges nyatakake yen
$\lforall[\Obj{v_0}][(\Obj{V^1_0}(\Obj v_0) \lor \lnot
  \Obj{V^1_0}(v_0))]$ bener kanggo saben cara menehi subhimpunan
saka~$\Domain{M}$ minangka nilai kanggo $\Obj{V^1_0}$,
utawa paling ora kanggo siji cara. Amarga saben himpunan mesthi
ngemot utawa ora ngemot sawijining objek tartamtu, loro-lorone
bener ing saben !!{structure}.
\end{document}
